% Phần Pipes của báo cáo — được \input từ main-vi.tex / main-en.tex.
\section{Bài toán 1: Pipes}

\subsection{Tìm hiểu đề bài và luật chơi}

Nhóm tìm hiểu Pipes qua việc chơi trực tiếp trên \url{https://www.puzzle-pipes.com/}
và quan sát hành vi của game gốc. Sau quá trình chơi thử, nhóm tóm tắt luật như sau:

\begin{itemize}
\item Bàn chơi $m\times m$ gồm các ô ống. Mỗi ô có \emph{hình dạng cố định}
(ống thẳng, khuỷu, chữ T, chữ thập, đầu cụt) và một \emph{định hướng} (hướng xoay)
có thể thay đổi; xoay là tác động duy nhất của người chơi.
\item Mỗi hình dạng có một tập cổng ở các phía Bắc--Đông--Nam--Tây. Xoay ô không
đổi số cổng, chỉ đổi hướng.
\item Mục tiêu: xoay sao cho \textbf{mọi} cổng mở đều gặp cổng mở tương ứng của ô
kề (không cổng nào chìa ra ngoài biên, không cổng nào \emph{dangling}), toàn bộ
các ô nối thành \textbf{một} mạng liên thông và mạng \textbf{không chứa chu trình}.
\item Nguồn nước (ô đánh dấu trên game gốc) chỉ để trực quan hóa: khi mạng hoàn
chỉnh, nước chảy từ nguồn tới mọi ô. Đổi vị trí nguồn không làm đổi tập các bàn
hợp lệ.
\end{itemize}

Nhìn theo phân loại của giáo trình, đây là bài toán \emph{một agent, tất định,
quan sát được toàn phần, rời rạc, tĩnh}. Đây là bài toán \emph{thỏa mãn}
(satisfaction): mọi bàn hợp lệ đều là lời giải, bài toán không yêu cầu tối ưu số
lần xoay. Trò chơi gốc cũng không tính điểm theo số bước xoay.

\subsection{Phát biểu bài toán tìm kiếm}

Gọi $n = m^2$, đánh số các ô theo thứ tự hàng--cột (row-major). Ô $i$ có hình dạng
cố định và miền định hướng $R_i$ gồm các xoay phân biệt của hình dạng đó. Nhóm
phát biểu bài toán như sau:

\begin{table}[h]
\centering\small
\begin{tabular}{@{}p{.20\textwidth}p{.74\textwidth}@{}}
\toprule
Thành phần & Định nghĩa trong cài đặt\\
\midrule
Input & Bàn chơi đã xáo trộn (hình dạng + định hướng lưu sẵn), kích thước, nguồn\\
Trạng thái & Tiền tố $s=(a_0,\dots,a_{k-1})$ các định hướng đã gán\\
Trạng thái đầu & Tiền tố rỗng $()$, biết trước toàn bộ hình dạng\\
Tác động & Gán một định hướng hợp lệ biên cho ô $k$ (ô kế tiếp theo row-major)\\
Chuyển tiếp & $s' = s + (a_k)$, chỉ nhận khi qua hết các kiểm tra khả thi\\
Kiểm tra đích & $k=n$ và bàn hoàn chỉnh thỏa mọi luật\\
Chi phí & 1 mỗi phép gán; không giải thuật nào dùng chi phí để xếp hạng\\
Lời giải & Một bộ định hướng đầy đủ hợp lệ\\
\bottomrule
\end{tabular}
\caption{Phát biểu bài toán tìm kiếm cho Pipes.}
\end{table}

Điểm cần nhấn: \emph{một bước của giải thuật} là gán \textbf{tuyệt đối} một định
hướng cho ô kế tiếp, khác với \emph{một bước của người chơi} là xoay 90 độ. Một
phép gán có thể tương ứng 0--3 lần xoay. Cả hai giải thuật dùng cùng mô hình gán,
nên việc so sánh chúng là công bằng. Nhóm cũng tách bạch trạng thái tìm kiếm
(tiền tố rỗng) với bàn xáo trộn lưu trong file input: bàn xáo trộn cung cấp miền
$R_i$ và thứ tự duyệt các xoay, còn tìm kiếm luôn bắt đầu từ tiền tố rỗng, không
dựa vào đáp án của generator.

\subsection{Không gian tìm kiếm và cắt tỉa dùng chung}

Số cách gán đầy đủ (chưa xét biên) là $U=\prod_i |R_i| \le 4^n$: ô khuỷu/T/đầu
cụt có 4 xoay phân biệt, ống thẳng có 2, ô chữ thập có 1. Với bàn 4$\times$4 điển
hình $U = 4^{16} \approx 4{,}29\times10^9$; với 7$\times$7 cận trên là
$4^{49} \approx 3{,}17\times10^{29}$. Cây gán có độ sâu $n$, hệ số phân nhánh tối
đa 4, và \textbf{mỗi tiền tố có đúng một cha} (thứ tự gán cố định), nên không cần
bảng closed để chống lặp trạng thái.

Trước khi tìm kiếm, các định hướng chìa cổng ra ngoài biên bị loại. Với mỗi ứng
viên con, hàm \code{compatible()} dùng chung cho cả hai giải thuật thực hiện ba
lớp cắt tỉa hợp lệ (sound):

\begin{enumerate}
\item \textbf{Khớp với láng giềng đã gán}: mỗi miền $R_i$ chưa gán bị thu hẹp chỉ
giữ định hướng khớp cổng của các ô đã gán (cả cổng có lẫn cổng \emph{không} có).
Miền nào rỗng thì chặn nhánh.
\item \textbf{Chặn chu trình}: nếu các ô đã gán tạo chu trình thì loại (dùng cấu
trúc disjoint-set).
\item \textbf{Kết nối khả dĩ}: dựng đồ thị \emph{lạc quan} giữa các ô — có cạnh
nếu hai ô \emph{có thể} mở cổng về phía nhau với một định hướng còn lại nào đó.
Nếu đồ thị này không liên thông thì không thể có mạng hoàn chỉnh, chặn nhánh.
\end{enumerate}

Tại gốc, kiểm tra nhanh $2(n-1)$ cổng: mạng cây trên $n$ đỉnh có đúng $n-1$ cạnh
$= 2(n-1)$ cổng; xoay không đổi số cổng nên tổng sai chứng tỏ vô nghiệm ngay từ
đầu. Điều ngược lại không đúng (đủ cổng chưa chắc có nghiệm). Cả ba lớp cắt tỉa
đều chỉ loại trạng thái \emph{không thể} dẫn tới nghiệm nên không mất tính đầy đủ.
DFS do đó là giải thuật mù về \emph{thứ tự duyệt}, nhưng vẫn là DFS có cắt tỉa
ràng buộc chứ không phải máy đếm xoay vét cạn.

\subsection{Giải thuật tìm kiếm mù: Depth-First Search}

DFS đi sâu theo nhánh con khả thi đầu tiên cho đến khi gặp đích hoặc bế tắc rồi
quay lui. Thứ tự duyệt các xoay là thứ tự lưu sẵn của từng ô — hoàn toàn không
dùng thông tin heuristic:

\begin{lstlisting}
SEARCH_DFS(input):
  build boundary-valid domains R_i
  if root feasibility fails: return UNSOLVABLE
  return VISIT(())

VISIT(s):
  if budget exhausted: return LIMIT_REACHED
  count one expansion
  if GOAL(s): return SOLVED(s)
  for a in boundary-valid rotations of tile |s|:
    child = s + (a)
    if FEASIBLE(child):           # 3 pruning layers
      r = VISIT(child)
      if r is SOLVED or LIMIT_REACHED: return r
  return UNSOLVABLE
\end{lstlisting}

Với không gian hữu hạn, cây có chu kỳ bị cắt bởi cấu trúc tiền tố, và cắt tỉa là
hợp lệ nên DFS \emph{đầy đủ} nếu đủ tài nguyên. Lời giải đầu tiên tìm được không
đảm bảo tối ưu số phép gán. Độ phức tạp thời gian xấu nhất là hàm mũ theo $n$;
về bộ nhớ, DFS giữ đường hiện hành và các nhánh chưa xét — cài đặt Python giữ bản
sao tiền tố và khung generator cho mỗi mức, khoảng $O(n^2)$ trên một nhánh.

\subsection{Giải thuật heuristic: Greedy Best-First}

Greedy luôn pop khỏi hàng ưu tiên (min-heap) trạng thái có ít \emph{lựa chọn còn
lại} nhất:
\[
h(s) = \sum_{i \text{ chưa gán}} \log_2 |D_i(s)|, \qquad
\text{priority}(s) = (h(s),\ -\text{depth}(s),\ \text{insertion order})
\]
với $D_i(s)$ là miền biên-hợp-lệ của ô $i$ sau khi khớp với các láng giềng đã
gán. Tổng là logarit của tích các miền còn lại: ví dụ miền $(1,2,4)$ cho
$h=0+1+2=3$, còn $(1,1,2)$ cho $h=1$ — trạng thái B ít ``tự do'' hơn và được ưu
tiên. Nhóm chọn heuristic này vì: (i) nó rẻ, chỉ cần miền đã tính trong bước cắt
tỉa; (ii) mọi lời giải ở cùng độ sâu đều cần đúng $n$ phép gán nên $n-k$ không
phân biệt được gì, còn tổng log-miền ưu tiên nhánh ít phân nhánh nhất. $h$ ở đây
\emph{không admissible} — nó đo độ bất định còn lại, không phải cận dưới chi phí;
và không dùng thành phần $g$.

\begin{lstlisting}
SEARCH_GREEDY(input):
  build boundary-valid domains; root feasibility check
  push () into min-priority queue with priority 0
  while queue is not empty:
    if budget exhausted: return LIMIT_REACHED
    s = pop minimum-priority state;  count one expansion
    if GOAL(s): return SOLVED(s)
    for a in boundary-valid rotations of tile |s|:
      child = s + (a);  D = FEASIBLE(child)
      if D exists: push child with (h(child), -|child|, serial)
  return UNSOLVABLE
\end{lstlisting}

Trên không gian tiền tố hữu hạn, vô chu kỳ, việc giữ toàn bộ ứng viên khả thi
làm Greedy đầy đủ nếu đủ tài nguyên; nhưng Greedy phải giữ cả frontier trong heap
nên bộ nhớ có thể tăng hàm mũ.

\subsection{Hiện thực và demo trực quan}

Cài đặt nằm trong package \code{pipes\_problem}: \code{model.py} (luật, validation,
sinh bàn), \code{search.py} (hai giải thuật, cắt tỉa, đo thời gian/bộ nhớ),
\code{app.py} (giao diện pygame), \code{worker.py} (process đo độc lập),
\code{benchmark.py} (chạy batch). Giao diện demo:

\begin{itemize}
\item Chơi thủ công: click xoay (Shift để ngược chiều), right-click khóa ô,
Undo/Reset; ống có nước hiển thị màu xanh chảy từ nguồn.
\item \textbf{Demo tìm kiếm từng bước}: chọn DFS hoặc Greedy rồi \emph{Run /
Pause / Next / Speed}. Mỗi sự kiện hiển thị đúng loại: \emph{expand} (thăm/pop
tiền tố), \emph{assign} (DFS chọn con), \emph{queue} (Greedy push), \emph{reject}
(định hướng bị cắt tỉa — ô viền đỏ kèm lý do), \emph{backtrack}. Nút \emph{Next}
cho phép phân biệt \emph{mở rộng trạng thái} với \emph{sinh successor}.
\item \emph{Show solution} phát lại riêng bộ gán cuối cùng; bàn trở về trạng thái
đầu trước khi chạy. Kết quả \emph{vô nghiệm} và \emph{chạm giới hạn} hiển thị
thành hai thông báo phân biệt.
\end{itemize}

\subsection{Tạo input}

Nhóm sinh input bằng generator của project (không chép tay từ web): 9 map mặc
định gồm 3 kích thước $4\times4, 5\times5, 7\times7$ $\times$ 3 seed 261--263,
định dạng JSON gọn (\code{size}, \code{source}, \code{tiles}), tên file là định
danh input. Ngoài ra có \code{maps/custom/cdai.json} — map được hiệu chỉnh tay
thành \emph{vô nghiệm} để kiểm tra thông báo và hành vi cạn không gian. Generator
runtime cũng cho sinh map theo \code{-{}-size}/\code{-{}-seed}. Đây là input
\emph{mẫu của nhóm} (không phải map chính thức của trang web) và generator không
lưu nghiệm vào file map.

\subsection{Kết quả thực nghiệm: thời gian và bộ nhớ}

\textbf{Giao thức đo.} Chạy trên máy Linux x86\_64 (kernel 6.8), Python 3.13.14,
pygame-ce 2.5.8. Cả hai giải thuật nhận \emph{cùng} 9 map, cùng thứ tự gán
row-major, cùng bộ cắt tỉa, giới hạn 30 giây và 200{,}000 lần mở rộng mỗi lượt.
Mỗi cặp (map, giải thuật) chạy 5 lần: thời gian là solve headless bằng
\code{perf\_counter} (không vẽ hình); bộ nhớ là đỉnh cấp phát Python đo bằng
\code{tracemalloc} trong một replay riêng, chỉ gộp khi hai lần cho cùng status,
cùng bàn kết quả và cùng số lần mở rộng. Bảng số liệu đầy đủ (90 lượt) lưu tại
\code{docs/benchmarks/2026-10-25-pipes.csv}; các bảng dưới là trung vị (median)
theo từng input.

\begin{table}[h]
\centering\small
\begin{tabular}{@{}llrrr@{}}
\toprule
Input & Giải thuật & Solved/lượt & Thời gian (ms) & Bộ nhớ (KiB)\\
\midrule
pipes-4x4-261 & DFS    & 5/5 & 3.380  & 26.836\\
              & Greedy & 5/5 & 2.866  & 11.695\\
\midrule
pipes-4x4-262 & DFS    & 5/5 & 3.047  & 27.078\\
              & Greedy & 5/5 & 3.386  & 12.359\\
\midrule
pipes-4x4-263 & DFS    & 5/5 & 3.149  & 26.984\\
              & Greedy & 5/5 & 2.741  & 12.500\\
\midrule
pipes-5x5-261 & DFS    & 5/5 & 5.352  & 53.203\\
              & Greedy & 5/5 & 7.094  & 19.500\\
\midrule
pipes-5x5-262 & DFS    & 5/5 & 5.906  & 53.344\\
              & Greedy & 5/5 & 5.882  & 19.219\\
\midrule
pipes-5x5-263 & DFS    & 5/5 & 4.564  & 53.297\\
              & Greedy & 5/5 & 5.712  & 19.219\\
\midrule
pipes-7x7-261 & DFS    & 5/5 & 23.247 & 145.852\\
              & Greedy & 5/5 & 27.049 & 27.594\\
\midrule
pipes-7x7-262 & DFS    & 5/5 & 21.938 & 145.930\\
              & Greedy & 5/5 & 26.669 & 27.070\\
\midrule
pipes-7x7-263 & DFS    & 5/5 & 40.785 & 145.852\\
              & Greedy & 5/5 & 38.246 & 29.094\\
\bottomrule
\end{tabular}
\caption{Pipes — thời gian và đỉnh bộ nhớ theo từng input (median 5 lượt, đo
25/10/2026). Cả 90/90 lượt đều \emph{solved}. Nguồn:
\code{docs/benchmarks/2026-10-25-pipes.csv}.}
\label{tab:pipes-per-input}
\end{table}

\begin{table}[h]
\centering\small
\begin{tabular}{@{}llrrr@{}}
\toprule
Kích thước & Giải thuật & Solved/lượt & Median time (ms) & Median memory (KiB)\\
\midrule
$4\times4$ & DFS    & 15/15 & 3.149  & 26.984\\
           & Greedy & 15/15 & 3.097  & 12.359\\
$5\times5$ & DFS    & 15/15 & 5.261  & 53.297\\
           & Greedy & 15/15 & 6.157  & 19.219\\
$7\times7$ & DFS    & 15/15 & 24.469 & 145.852\\
           & Greedy & 15/15 & 31.695 & 27.594\\
\bottomrule
\end{tabular}
\caption{Pipes — tổng hợp theo kích thước (3 map $\times$ 5 lượt mỗi ô số liệu).}
\label{tab:pipes-summary}
\end{table}

\textbf{Giải thích kết quả.}

\begin{itemize}
\item \textbf{Thời gian tăng nhanh theo kích thước} (3--5 ms ở 4$\times$4 lên
22--41 ms ở 7$\times$7): số biến tăng từ 16 lên 49 và chi phí kiểm tra khả thi
mỗi successor tăng theo số láng giềng đã gán, đúng dự đoán hàm mũ về không gian
tìm kiếm.
\item \textbf{DFS nhanh hơn Greedy ở 5$\times$5 và 2/3 map 7$\times$7.} Trên các
input dễ, cắt tỉa đã khiến nhánh đầu tiên của DFS gần như thẳng tới nghiệm; hàng
ưu tiên của Greedy lúc đó chỉ tạo thêm chi phí push/pop heap và tính $h$ mà không
đổi thứ tự duyệt đáng kể. Đây là minh chứng cho việc heuristic \emph{không miễn
phí}: chi phí cài đặt phải được bù bằng chất lượng xếp hạng.
\item \textbf{Trên map khó nhất (pipes-7x7-263), Greedy nhanh hơn DFS}
(38.2 so với 40.8 ms): map này buộc tìm kiếm phải lui vài nhánh thật; xếp hạng
theo ``ít lựa chọn còn lại'' giúp nhánh gấp gáp (miền thu hẹp nhanh) được xét
trước, tiết kiệm bớt phần dò thừa. Chênh lệch nhỏ (5\%) nhưng xuất hiện ở đúng
loại input mà heuristic có giá trị.
\item \textbf{Bộ nhớ: Greedy thấp hơn DFS 4--5 lần} (12--29 KiB so với
27--146 KiB). Nghịch với suy đoán lý thuyết ``frontier của Greedy có thể hàm mũ'',
số đo tracemalloc phản ánh \emph{cấp phát Python thực tế của cài đặt}: DFS giữ
bản sao tiền tố cho từng mức đệ quy và khung generator treo (mỗi mức một tuple
copy), trong khi Greedy trên input dễ chỉ giữ vài chục ứng viên do cắt tỉa gạt
gần hết. Trên input khó hơn nhiều, frontier của Greedy vẫn có nguy cơ phình hàm
mũ — cần thí nghiệm bổ sung để khẳng định.
\item \textbf{Cả 90 lượt đều solved} trong giới hạn; không có run nào chạm 30 giây
hay 200k expansions trên bộ input này.
\end{itemize}

\subsection{Hạn chế}

Bộ input mặc định khá dễ với cắt tỉa hiện tại nên hai giải thuật chưa bị ép tới
giới hạn; kết luận về ``kẻ thắng'' chỉ có giá trị trên loại input này. Nhóm chưa
đo riêng đóng góp từng lớp cắt tỉa (ablation), chưa thêm đủ nhiều map vô nghiệm,
và chưa thử heuristic MRV. Các hạn chế này được ghi nhận trong
\code{docs/pipes-analysis.md} và là việc làm tiếp theo.
